\section{Ethics, Safety, and Method Limitations}
\label{sec:study_design_ethics_safety_limitations}

The study involves physical exertion on an ergometer while wearing a VR headset. The procedure must therefore include basic safety screening, clear stop criteria, and the option to withdraw at any time without giving a reason.

Participants will be informed that they may experience physical fatigue, discomfort, dizziness, or VR-related sickness symptoms. They will be instructed to stop immediately if they feel unsafe or unwell. The experimenter will monitor the session and pause or terminate the trial if necessary.

\drafttodo{TODO:CONFIRM consent form, data protection wording, local ethics or supervisor approval requirements, and final safety checklist.}

The first study version deliberately avoids threat stimuli, physiological response measures, and hand-offset probes. This reduces additional ethical complexity and keeps the first evaluation focused on the main representation comparison.

The method has expected limitations. It relies mainly on self-report, uses a prototype-specific implementation, and may involve a small or convenience-based sample. These limitations mean that the results should be interpreted as evidence about the studied XR rowing setup rather than as general conclusions about all VR rowing systems or all embodiment mechanisms.
